\section{Question Analysis}

\subsection{Data Analysis}
The raw table carries columns beyond the model's needs, so the target variable is first
located and matched to the physical quantity the model expects. [Then list each
operation and its reason: fill / truncate / interpolate.]

\subsection{Analysis of Problem 1}
[Known, unknown, framework, output — four beats in one paragraph.]

\subsection{Analysis of Problem 2}
[Same four beats; name the key difficulty in one sentence.]

\subsection{Analysis of Problem 3}
[Same four beats; state what information is missing.]

% 【Section writing prompt】
% 2.1 每个处理都要给"为什么"，不是报流水账：只写"缺失值填 0"没人信，写清
% "填 0 表示无输入，是保守策略，避免注入不合理扰动"才有说服力。
% 2.2–2.4 一段一问：已知 → 未知 → 框架 → 输出，再加一句"关键难点在哪"。
% 能从句面直接抄回题面的句子，删掉。
